\chapter{Đặt vấn đề}
\label{ch:dat-van-de}

\section{Bối cảnh}

Ở trường phổ thông, một bài kiểm tra thường kết thúc khi giáo viên trả điểm. Học sinh nhận lại bài,
nhìn những câu bị gạch chéo, rồi phần lớn dừng ở đó. Thứ thực sự làm thay đổi kết quả học tập là
việc hiểu ra mình sai ở đâu và làm lại được, nhưng việc đó hiếm khi xảy ra, vì nó cần một cuộc trao
đổi riêng cho từng em về từng câu.

Giáo viên không phải không muốn làm. Vấn đề nằm ở khối lượng công việc. Một lớp bốn mươi học sinh
làm một đề mười câu sẽ sinh ra bốn trăm lượt trả lời, và mỗi lượt sai có thể sai vì một lý do khác
nhau. Giải thích riêng cho từng em, rồi ra thêm một câu tương tự để kiểm xem em đã sửa được chưa, là
công việc tăng theo sĩ số và không có cách nào rút gọn.

Kriky là hệ thống được xây để gánh phần công việc đó, nhưng không lấy đi quyền quyết định của giáo
viên. Hình~\ref{fig:domain-context} cho thấy ranh giới của hệ thống: hai actor ở hai phía, cùng
những gì đi vào và đi ra.

\begin{landscape}
\begin{figure}[p]
  \centering
  \includegraphics[width=\linewidth,height=\textheight,keepaspectratio]{domain-context.pdf}
  \caption[Domain Context]{Domain Context --- ranh giới hệ thống, hai actor, và những gì đi qua
  ranh giới ấy}
  \label{fig:domain-context}
\end{figure}
\end{landscape}

\section{Vấn đề cần giải quyết}

Phần tốn thời gian nhất của giáo viên không phải là gõ ra đề bài, mà là năm việc sau. Cả năm đều là
công việc chuyên môn:

\begin{itemize}
  \item Soạn đề bám đúng mục tiêu học tập, thay vì gom câu hỏi cho đủ số lượng.
  \item Tạo phương án nhiễu có ý nghĩa. Một phương án sai ngẫu nhiên không nói lên điều gì, còn
        một phương án nhiễu tốt là phương án mà học sinh mắc một lỗi cụ thể sẽ chọn.
  \item Chấm bài và viết nhận xét cho từng học sinh.
  \item Phân định học sinh sai vì thiếu kiến thức, vì suy luận hỏng, vì tính nhầm, hay vì hiểu
        sai một khái niệm.
  \item Ra câu biến thể của đúng câu học sinh vừa làm sai, để kiểm xem em đã sửa được chỗ sai đó
        hay chỉ nhớ được đáp án.
\end{itemize}

Phía học sinh, nhu cầu ngược lại nhưng khớp vào: phản hồi cần đến nhanh, và bài luyện tập cần đúng
với điểm yếu của riêng mình. Một phản hồi đúng nhưng đến sau hai tuần thì không còn tác dụng, vì học
sinh đã quên mình nghĩ gì lúc làm bài.

\section{Mục tiêu sản phẩm}

Hệ thống hướng tới bốn mục tiêu nghiệp vụ:

\begin{enumerate}
  \item Hỗ trợ giáo viên tạo đề từ một yêu cầu bằng ngôn ngữ tự nhiên. Một đề có hai nguồn câu
        hỏi: nội dung sinh mới, và kho câu hỏi đã dùng qua.
  \item Hỗ trợ chấm bài và sinh nhận xét có căn cứ. Mỗi nhận xét phải dẫn được về một bằng
        chứng, thay vì là một lời khuyên chung chung.
  \item Phân tích lỗi sai hoặc hiểu nhầm khái niệm của học sinh.
  \item Bắt buộc học sinh chữa lại những câu đã sai, ngay trong cùng bài kiểm tra.
\end{enumerate}

Mục tiêu thứ tư là thứ phân biệt hệ thống này với một công cụ sinh đề thông thường.
Chương~\ref{ch:y-tuong} dành phần lớn dung lượng cho nó.

\section{Phạm vi}

Hệ thống dựng trọn một vòng nghiệp vụ, thay vì dựng sâu một khâu:

\begin{itemize}
  \item Giáo viên tạo, xem lại và phát hành đề. Việc duyệt bao gồm cả lời giải và ánh xạ từ
        phương án nhiễu sang lỗi, không chỉ câu hỏi và đáp án.
  \item Giai đoạn 2: Học sinh làm bài và nộp bài.
  \item Hệ thống chấm bài và tra lỗi sai từ phương án nhiễu mà học sinh đã chọn.
  \item Kết quả có độ tin cậy thấp, hoặc có dấu hiệu bất thường, đi vào hàng đợi chờ giáo viên
        xem lại.
  \item Giai đoạn 2: hệ thống giải thích lỗi và sinh câu biến thể, học sinh làm lại cho tới khi điểm
        của từng câu được chốt.
\end{itemize}
